\chapter{Portfolio Construction I}\label{ch:rs:portfolio-construction-i}

An unconstrained mean--variance optimiser is given a hundred large stocks, an honest alpha forecast and the sample covariance of their last six months of daily
returns. It builds a long--short book whose gross exposure averages sixty-one times its capital, forecasts it at 36 per cent annual volatility and an
information ratio of 4.25, and rebuilds it every month for eight years. The book realises a volatility of 207 per cent and an information ratio of 0.92 before
costs; it trades eighty-four times its capital a month, and after costs of ten basis points a trade its information ratio is 0.43. The same forecasts, through
chapter~24's risk model and the constraints a real book carries, realise 1.77 after costs. The optimiser did what it was told: it found the directions in which
the inputs were most wrong. This chapter builds \texttt{firm.portcons}, a construction layer over Book~4's \texttt{firm.portopt}, adds constraints one at a time,
and prices each one with its shadow price.

\section{Mean--variance with an alpha and a risk model}

\begin{definition}[Mean--variance optimisation, risk-aversion parameter, efficient frontier]\label{def:rs:portfolio-construction-i:mv}
\emph{Mean--variance optimisation}\index{mean--variance optimisation} chooses the weights $w$ that maximise $\alpha^\top w - \frac{\gamma}{2} w^\top\Sigma w$
subject to constraints, for an alpha forecast $\alpha$ and a covariance forecast $\Sigma$ over the same horizon (Markowitz, 1952). The
\emph{risk-aversion parameter}\index{risk-aversion parameter} $\gamma$ sets the trade between expected return and variance. The \emph{efficient
frontier}\index{efficient frontier} is the set of portfolios, over all $\gamma$, with the highest expected return for their risk.
\end{definition}

\begin{proposition}[The unconstrained and the equality-constrained optimum]\label{prop:rs:portfolio-construction-i:mv}
Without constraints the optimum is $w^* = \Sigma^{-1}\alpha/\gamma$, with ex-ante information ratio $\sqrt{\alpha^\top\Sigma^{-1}\alpha}$ whatever $\gamma$. With
equality constraints $A^\top w = 0$ it is $w^* = \Sigma^{-1}(\alpha - A\lambda)/\gamma$, where $\lambda$, the shadow prices, solve $A^\top\Sigma^{-1}(\alpha - A\lambda) = 0$.
\end{proposition}

\begin{proof}
Setting the gradient of $\alpha^\top w - \frac\gamma2 w^\top\Sigma w - \lambda^\top A^\top w$ to zero gives $\gamma\Sigma w = \alpha - A\lambda$; the constraint fixes
$\lambda$. At $w^* = \Sigma^{-1}\alpha/\gamma$, $\alpha^\top w^*/\sqrt{w^{*\top}\Sigma w^*} = \sqrt{\alpha^\top\Sigma^{-1}\alpha}$.
\end{proof}

The chapter's book holds the hundred largest names of \texttt{firm.synthmkt}, rebuilt at 95 month-ends from year~3 to year~10 and held over the following month's
days. Its alpha is the simulation's stock-specific expected return for the month (the persistent drift and the post-earnings drift), plus noise twice as large,
standardised and scaled by chapter~15's rule, $\alpha_i = \mathrm{IC}\,\sigma_i z_i$ with an \omterm{def:rs:anatomy-of-a-predictor:ic}{information coefficient} of 0.05 and $\sigma_i$ the name's residual
volatility; the noise is new every month. The risk model is chapter~24's at a monthly horizon; $\gamma = 12$; trades cost ten basis points of the weight traded.
The alpha is the truth plus noise on purpose: the chapter studies what construction does to a given forecast, not the forecast.

\section{The error maximiser}

\begin{definition}[Error maximisation]\label{def:rs:portfolio-construction-i:error}
\emph{Error maximisation}\index{error maximisation} is the tendency of an optimiser to load the directions in which its inputs are most in error: assets with
overestimated returns, underestimated variances or understated correlations, which look like the best trades because they are the worst estimates (Michaud,
1989, who called mean--variance optimisers ``estimation-error maximizers'').
\end{definition}

The opening's optimiser uses the forecast in return units (the same scale for every name, not divided by risk) and the sample covariance of 126 days for 100
names. With barely more days than names, the smallest eigenvalues of the sample covariance are far below the truth, and $\Sigma^{-1}$ amplifies them: the book
takes enormous offsetting positions in combinations of stocks that happened to move together in the last six months. Its gross exposure averages 61, its
forecast volatility 36\% against 207\% realised, and the two largest positions hold only 9\% of the gross on average (16\% at most): the errors are spread across
hedged combinations rather than concentrated in a pair of names. The ex-ante information ratio of 4.25 is the optimiser's measure of its own errors. Replacing
the inputs, not constraining the output, is the first fix: alphas scaled by the risk the optimiser sees (chapter~15), and a factor risk model whose few
parameters are estimated on many stocks (chapter~24). With nothing else changed but \omterm{def:rs:portfolio-construction-i:constraints}{dollar neutrality}, the book's forecast volatility becomes 13.7\% against
14.0\% realised, and its information ratio 1.39 after costs.

\section{Realistic constraints}

\begin{definition}[Dollar, beta and factor neutrality; name limit; liquidity constraint; 130/30 portfolio]\label{def:rs:portfolio-construction-i:constraints}
A book has \emph{dollar neutrality}\index{dollar neutrality} when its long and short weights sum to zero, \emph{beta neutrality}\index{beta neutrality} when its
\omterm{def:rs:performance-measurement:beta}{market beta} is zero, and satisfies a \emph{factor-neutrality constraint}\index{factor-neutrality constraint} when its exposure to a risk-model factor is zero. A
\emph{name limit}\index{name limit} bounds each position's weight; a \emph{liquidity constraint}\index{liquidity constraint} bounds it by a fraction of the
stock's average daily traded value relative to the book's capital. A \emph{130/30 portfolio}\index{130/30 portfolio} is long 130\% and short 30\% of its capital:
a long-only mandate given limited room to short.
\end{definition}

Constraints enter in the order a desk meets them. Each is a line in \texttt{rs\_portcons.build} (\cref{lst:rs:portfolio-construction-i:build}), and
\texttt{firm.portcons} turns the set into one quadratic programme for \texttt{firm.portopt.qp}:

\begin{center}
\small
\begin{tabular}{@{}lccccccc@{}}
\toprule
 & \multicolumn{3}{c}{information ratio} & & & & forecast / \\
\cmidrule(lr){2-4}
cumulative setting & ex ante & realised & net & turnover & gross & TC & realised vol.\\
\midrule
naive (sample covariance) & 4.25 & 0.92 & 0.43 & 84 & 61 & 0.44 & 36\% / 207\%\\
dollar neutral & 1.64 & 1.89 & 1.39 & 5.80 & 4.44 & 0.97 & 13.7\% / 14.0\%\\
+ beta, industry, style neutral & 1.59 & 1.72 & 1.22 & 5.70 & 4.35 & 0.92 & 13.3\% / 13.6\%\\
+ \omterm{def:rs:portfolio-construction-i:constraints}{name limits} of 3\% & 1.41 & 1.29 & 0.83 & 2.50 & 2.21 & 0.82 & 6.4\% / 6.6\%\\
+ gross at most 2 & 1.44 & 1.40 & 0.94 & 2.37 & 2.00 & 0.83 & 6.1\% / 6.3\%\\
+ liquidity, 5\% of daily value & 1.23 & 0.87 & 0.46 & 1.55 & 1.34 & 0.71 & 4.4\% / 4.4\%\\
+ turnover at most 0.5 a month & 1.00 & 1.94 & 1.77 & 0.53 & 1.01 & 0.58 & 3.5\% / 3.6\%\\
\bottomrule
\end{tabular}
\end{center}

Turnover is one-way monthly, as a multiple of capital; TC is the \omterm{def:rs:from-signal-to-forecast:breadth}{transfer coefficient} (chapter~15); the liquidity limit assumes \$1 billion of capital. Each
constraint lowers the ex-ante information ratio, by construction: the optimiser loses room. Realised information ratios over 7.9 years have standard errors of
about 0.5, so most of the differences between neighbouring rows are noise; the robust facts are the first row, the effect of costs, and the last row. Neutrality
to the risk model's factors costs little ex ante (1.64 to 1.59) because the alpha is stock-specific; Jacobs, Levy and Starer showed that dollar and \omterm{def:rs:portfolio-construction-i:constraints}{beta neutrality}
are in general suboptimal, and here they are nearly free. \omterm{def:rs:portfolio-construction-i:constraints}{Name limits} and the gross limit halve the book's scale and its costs. The liquidity limit binds on 28 of the hundred names on average and lowers the ex-ante information ratio from 1.44 to 1.23. Jagannathan and Ma explained why constraints that are wrong on paper help in
practice: they act as shrinkage on estimation error, the error maximiser's fuel.

\begin{omfigure}
\begin{tikzpicture}
\begin{axis}[width=12cm,height=5.6cm,ybar,bar width=6pt,ylabel={information ratio},xtick={0,1,2,3,4,5,6},
  xticklabels={naive,dollar,factor,names,gross,liquidity,turnover},tick label style={font=\small},label style={font=\small},
  xmin=-0.6,xmax=6.6,ymin=0,ymax=4.5,grid=major,grid style={gray!25},legend columns=3,
  legend style={font=\footnotesize,at={(0.5,-0.2)},anchor=north,draw=none,fill=none,/tikz/every even column/.append style={column sep=8pt}},
  table/col sep=comma]
\addplot[fill=gray!40,draw=gray] table[x=level,y=exante]{figdata/research/25-portfolio-construction-i/ir.csv};
\addplot[fill=omDef!40,draw=omDef] table[x=level,y=realised]{figdata/research/25-portfolio-construction-i/ir.csv};
\addplot[fill=omDef!85,draw=omDef] table[x=level,y=net]{figdata/research/25-portfolio-construction-i/ir.csv};
\legend{ex ante,realised,realised after costs}
\end{axis}
\end{tikzpicture}

\omcaption{Information ratios of the hundred-name book as constraints are added (each label names the constraint added to those on its left), 95 monthly rebalances. Ex-ante ratios fall with every constraint; realised ones
do not follow, and the turnover limit gives the best book after costs. Data: \texttt{rs\_portcons.summary}.}\label{fig:rs:portfolio-construction-i:ir}
\end{omfigure}

\section{Turnover control}

\begin{definition}[Turnover penalty]\label{def:rs:portfolio-construction-i:turnover}
A \emph{turnover penalty}\index{turnover penalty} subtracts $\kappa\,\lVert w - w_0\rVert_1$ from the objective, a linear cost on the weight traded from the
current book $w_0$; a turnover limit bounds $\lVert w - w_0\rVert_1$ instead.
\end{definition}

The turnover limit is the constraint that earns. It cuts the book's monthly trading from 1.55 to 0.53 times capital and its costs from 1.9\% to 0.6\% a year; its
realised information ratio rises from 0.87 to 1.94 before costs, which costs cannot explain. The reason is in the forecast: the planted drift persists for years,
while the noise added to it is new every month. A book that may move only part of the way towards each month's target averages several months' forecasts, and
averaging removes noise. The \omterm{def:rs:from-signal-to-forecast:breadth}{transfer coefficient}, the correlation between the book's risk-adjusted weights and the forecast (Clarke, de Silva and Thorley), falls
from 0.97 to 0.58, and it is right to fall: the forecast it transfers is noisier than the book. A \omterm{def:rs:portfolio-construction-i:turnover}{turnover penalty} or limit is therefore two things at once, a cost
model and a smoother of forecasts; chapter~26's multi-period rule makes the second explicit, and chapter~27 prices the first properly.

\begin{omfigure}
\begin{tikzpicture}
\begin{axis}[width=11cm,height=5.8cm,xlabel={years from the first rebalance},ylabel={cumulative return at 10\% volatility},tick label style={font=\small},
  label style={font=\small},xmin=0,xmax=8,ymin=-0.1,ymax=1.45,grid=major,grid style={gray!25},legend pos=north west,legend style={font=\footnotesize},
  yticklabel style={/pgf/number format/fixed,/pgf/number format/precision=1},scaled y ticks=false,table/col sep=comma]
\addplot[black!50,thick,dashed] table[x=year,y=naive]{figdata/research/25-portfolio-construction-i/cum.csv};
\addplot[omDef!60,thick] table[x=year,y=dollar]{figdata/research/25-portfolio-construction-i/cum.csv};
\addplot[black!70,thick] table[x=year,y=gross]{figdata/research/25-portfolio-construction-i/cum.csv};
\addplot[omDef,very thick] table[x=year,y=turnover]{figdata/research/25-portfolio-construction-i/cum.csv};
\legend{naive,dollar neutral,+ gross limit,+ turnover limit}
\end{axis}
\end{tikzpicture}

\omcaption{Cumulative returns after costs, each book scaled to 10\% annual volatility so that the slopes compare information ratios. Data:
\texttt{rs\_portcons.run}.}\label{fig:rs:portfolio-construction-i:cum}
\end{omfigure}

\section{Shadow prices and what each constraint costs}

Every constraint has a shadow price, the rate at which the objective would improve if it were relaxed (Book~4, chapter~23). \texttt{firm.portcons} returns it, and
more: the vector through which each constraint bends the optimum.

\begin{proposition}[Decomposing the information ratio by constraint]\label{prop:rs:portfolio-construction-i:decomp}
At the optimum of a constrained mean--variance problem, $\alpha - \gamma\Sigma w^* = \sum_k g_k$, where $g_k$ collects constraint $k$'s rows weighted by their
multipliers. The unconstrained ex-ante $\mathrm{IR}^2 = \alpha^\top\Sigma^{-1}\alpha$ exceeds $\gamma^2 w^{*\top}\Sigma w^*$ by exactly $\sum_k \delta_k$, with
$\delta_k = g_k^\top\Sigma^{-1}(2\alpha - G)$ and $G = \sum_k g_k$.
\end{proposition}

\begin{proof}
The first statement is the stationarity condition of the Karush--Kuhn--Tucker system restricted to the weights (split variables for absolute values contribute
through the rows that link them to $w$). Then $\gamma^2 w^{*\top}\Sigma w^* = (\alpha - G)^\top\Sigma^{-1}(\alpha - G) = \alpha^\top\Sigma^{-1}\alpha - \sum_k
g_k^\top\Sigma^{-1}(2\alpha - G)$.
\end{proof}

The quantity $\gamma^2 w^\top\Sigma w$ is the squared information ratio of the alpha the constraints leave to the optimiser, and $\delta_k$ is the share of the
forecast's information that constraint $k$ takes away; the shares add up (\cref{lst:rs:portfolio-construction-i:decomp}). For the fully constrained book,
averaged over the 95 months, the turnover limit takes 44.0\% of the unconstrained $\mathrm{IR}^2$, liquidity 39.4\%, the neutralities 7.0\%, the \omterm{def:rs:portfolio-construction-i:constraints}{name limits} 3.1\%;
the gross limit and \omterm{def:rs:portfolio-construction-i:constraints}{dollar neutrality} take nothing, because once liquidity and turnover have shrunk the book its gross is 1.0 and never binds. On an average month
28 liquidity bounds bind, 2.5 \omterm{def:rs:portfolio-construction-i:constraints}{name limits}, and the turnover limit always. Ex ante, turnover and liquidity cost most; realised, the turnover limit earned more than it
cost (\cref{fig:rs:portfolio-construction-i:decomp}). That is the named result, and the reason to report both: shadow prices say what a constraint costs if the
forecast is right, and a \omterm{def:rs:vectorised-backtests:backtest}{backtest} after costs says whether it was.

\begin{omfigure}
\begin{tikzpicture}
\begin{axis}[xbar,width=9.5cm,height=4.8cm,bar width=9pt,xlabel={share of the unconstrained ex-ante $\mathrm{IR}^2$},symbolic y coords={factor neutral,
  name limits,gross limit,liquidity,turnover limit},ytick=data,tick label style={font=\small},label style={font=\small},xmin=0,xmax=0.5,
  grid=major,grid style={gray!25},xticklabel style={/pgf/number format/fixed,/pgf/number format/precision=1},scaled x ticks=false,table/col sep=comma]
\addplot[fill=omDef!55,draw=omDef] table[y=group,x=share]{figdata/research/25-portfolio-construction-i/decomp.csv};
\end{axis}
\end{tikzpicture}

\omcaption{The shadow-price decomposition of the fully constrained book: the average share of the unconstrained squared ex-ante information ratio each constraint
takes away. Data: \texttt{rs\_portcons.decomposition}.}\label{fig:rs:portfolio-construction-i:decomp}
\end{omfigure}

\section{Solving at scale}

\texttt{firm.portcons} writes the problem in the factor form of the risk model, $\Sigma = XFX^\top + D$, and adds split variables only for the absolute values
that the gross and turnover terms need, so a hundred names become at most five hundred variables. \texttt{firm.portopt.qp} is a dense interior-point solver, whose cost
grows with the cube of the number of variables: enough for a hundred names and 95 months in under a minute, not for a firm book of three thousand names rebuilt
daily, which needs the structure exploited. Introducing the \omterm{def:rs:risk-models:families}{factor exposures}
$y = X^\top w$ as variables makes the quadratic term $y^\top F y + w^\top D w$, diagonal plus a small dense block, and a sparse solver or an operator-splitting
method (Book~4's ADMM) then scales with the number of names rather than its cube. Warm starts from last month's solution, and screening names whose alpha
cannot pay for their cost, do the rest.

\section{Tutorial: the sixty-one-times book}\label{tut:rs:portfolio-construction-i:tut}

\begin{tutorial}
\textbf{Goal.} Build the hundred-name book with seven cumulative sets of constraints, 95 months, and record ex-ante and realised information ratios, turnover,
\omterm{def:rs:from-signal-to-forecast:breadth}{transfer coefficient}, shadow prices and the decomposition. \textbf{End state:} the table, \cref{fig:rs:portfolio-construction-i:ir,fig:rs:portfolio-construction-i:cum,fig:rs:portfolio-construction-i:decomp}.
\begin{tutsteps}
\item \textbf{The problem}, one constraint set at a time.
\omcode{code/research/25-portfolio-construction-i/python/rs_portcons.py}{79}{106}{The chapter's constraint sets.}\label{lst:rs:portfolio-construction-i:build}
\item \textbf{Shadow-price vectors and the decomposition.}
\omcode{code/firm/portcons/firm_portcons.py}{147}{170}{Each constraint's pull on the optimum, and its share of the information ratio.}\label{lst:rs:portfolio-construction-i:decomp}
\item \textbf{Run} \texttt{rs\_portcons.summary(k)} for $k = 0, \dots, 6$, \texttt{decomposition(6)} and \texttt{fig\_portcons.py}.
\end{tutsteps}
\textbf{What to change next.} Replace the turnover limit by a penalty and trace net information ratio against $\kappa$; make the noise persistent (the same
draw each month) and see the turnover limit lose its advantage.
\end{tutorial}

\section{Build: the construction layer}\label{bld:rs:portfolio-construction-i:portcons}

\begin{build}
\bfield{Purpose} Every book the firm trades is built by one layer: the alpha, the risk model and the costs in, the constraints named, the trades and the price of
each constraint out.
\bfield{Interface} \texttt{Problem(alpha, X, F, spec, gamma, w0)}; \texttt{.equality(row, value, name)}, \texttt{.neutral\_factors(cols, names)},
\texttt{.bounds(lo, hi, name)}, \texttt{.liquidity(adv, participation, capital, name)}, \texttt{.gross(limit, name)}, \texttt{.turnover(limit, name)},
\texttt{.turnover\_penalty(kappa)}; \texttt{.solve()} returns weights, ex-ante information ratio, shadow prices, binding counts and pull vectors;
\texttt{.ir\_decomposition(result)}; \texttt{ir\_ex\_ante}; \texttt{trade\_list(w, w0, capital, prices, names)}.
\bfield{Rules} Alphas are scaled by the risk the optimiser sees; the covariance is a factor model, never a sample covariance of more names than it has periods;
every constraint is named, and its shadow price and share are reported with the book.
\bfield{Acceptance tests} \texttt{code/firm/portcons/tests/}: the closed form without constraints; neutralities met, and multipliers equal to the KKT ones; bounds,
gross and turnover limits met and binding; the gross limit's shadow price equal to the objective's finite difference; a large \omterm{def:rs:portfolio-construction-i:turnover}{turnover penalty} holds the book still;
liquidity bounds; pull vectors summing to $\alpha - \gamma\Sigma w$ and a decomposition that adds up; a trade list in shares.
\bfield{Stretch} A factor-form solver for thousands of names; the eigenfactor adjustment of chapter~24's model inside the optimiser; tax lots.
\end{build}

\begin{omsources}
\item H.~Markowitz, ``Portfolio selection'', \emph{Journal of Finance} 7(1), 1952.
\item R.~O.~Michaud, ``The Markowitz optimization enigma: is `optimized' optimal?'', \emph{Financial Analysts Journal} 45(1), 1989.
\item R.~Jagannathan and T.~Ma, ``Risk reduction in large portfolios: why imposing the wrong constraints helps'', \emph{Journal of Finance} 58(4), 2003.
\item B.~I.~Jacobs, K.~N.~Levy and D.~Starer, ``On the optimality of long--short strategies'', \emph{Financial Analysts Journal} 54(2), 1998.
\item R.~Clarke, H.~de Silva and S.~Thorley, ``Portfolio constraints and the fundamental law of active management'', \emph{Financial Analysts Journal} 58(5), 2002.
\end{omsources}

\section{Exercises}

\begin{exercise}[$\star$]\label{exo:rs:portfolio-construction-i:1}
Two uncorrelated assets have alphas of 2\% and 1\% and volatilities of 20\% and 10\%; $\gamma = 5$. Give the unconstrained weights and the ex-ante information ratio.
\end{exercise}

\begin{exercise}[$\star$]\label{exo:rs:portfolio-construction-i:2}
The naive book trades 84 times its capital a month at ten basis points. What do its costs come to in a year?
\end{exercise}

\begin{exercise}[$\star$]\label{exo:rs:portfolio-construction-i:3}
Give the standard errors of realised information ratios of 1.77 and 1.39 measured over 7.9 years.
\end{exercise}

\begin{exercise}[$\star\star$]\label{exo:rs:portfolio-construction-i:4}
Why does the sample covariance of a hundred stocks over 126 days lead to gross exposures of sixty times capital?
\end{exercise}

\begin{exercise}[$\star\star$]\label{exo:rs:portfolio-construction-i:5}
Explain how a turnover limit can raise the realised information ratio before costs.
\end{exercise}

\begin{exercise}[$\star\star$]\label{exo:rs:portfolio-construction-i:6}
With \$1 billion of capital and positions limited to 5\% of a stock's average daily traded value, what is the largest weight in a stock trading \$10 million a
day?
\end{exercise}

\begin{exercise}[$\star\star\star$]\label{exo:rs:portfolio-construction-i:7}
\emph{Coding.} With \texttt{rs\_portcons.decomposition(6)} and \texttt{run(6)}, explain why the gross limit's share is zero although it is in the problem.
\end{exercise}

\begin{exercise}[$\star\star\star$]\label{exo:rs:portfolio-construction-i:8}
\emph{Find the flaw.} ``The optimiser's ex-ante information ratio is 4.25 on a sample covariance of the last six months, so we will size the fund on it.''
\end{exercise}

\section{Problem: The Sixty-One-Times Book}

\begin{problem}[{Weekend problem --- constraints, priced}]\label{pb:rs:portfolio-construction-i:1}
The chapter's hundred-name book on \texttt{firm.synthmkt}.

\medskip\noindent\textbf{Part I --- The optimiser.}
\begin{enumerate}
\item Write the unconstrained optimum and its ex-ante information ratio.
\item How are the alpha and the risk model built, and why is the alpha the truth plus noise?
\item What does the naive book look like: gross, forecast and realised volatility, ex-ante and realised information ratios, turnover?
\item Why are its errors spread over hedged combinations rather than two names?
\end{enumerate}

\medskip\noindent\textbf{Part II --- Constraints.}
\begin{enumerate}[resume]
\item What does the factor risk model with \omterm{def:rs:portfolio-construction-i:constraints}{dollar neutrality} change?
\item What do factor neutralities cost ex ante, and why so little?
\item What do name and gross limits do to scale and costs?
\item How many liquidity bounds bind, and what do they cost ex ante?
\item Why are most differences in realised information ratio between neighbouring settings noise?
\end{enumerate}

\medskip\noindent\textbf{Part III --- Turnover.}
\begin{enumerate}[resume]
\item What do turnover and costs become with the turnover limit?
\item Why does the realised information ratio rise before costs?
\item Why does the \omterm{def:rs:from-signal-to-forecast:breadth}{transfer coefficient} fall, and is that bad?
\end{enumerate}

\medskip\noindent\textbf{Part IV --- The verdict.}
\begin{enumerate}[resume]
\item State the proposition behind the decomposition.
\item State the \emph{named result}: the share of the unconstrained ex-ante information each constraint takes, and what each earned after costs.
\item Why do the gross limit and \omterm{def:rs:portfolio-construction-i:constraints}{dollar neutrality} take nothing?
\item Which book would you trade, and what would you report with it?
\item When is a sample covariance acceptable in an optimiser?
\item What would Jagannathan and Ma say about the \omterm{def:rs:portfolio-construction-i:constraints}{name limits}?
\item How would you solve the problem for three thousand names?
\item In one sentence: what is an optimiser's ex-ante information ratio worth?
\end{enumerate}
\end{problem}

\section{Interview questions}

\begin{interviewq}[$\star$ \iqroles{researcher, trader}]\label{iq:rs:portfolio-construction-i:1}
Why do unconstrained mean--variance portfolios perform badly out of sample?
\end{interviewq}

\begin{interviewq}[$\star\star$ \iqroles{researcher}]\label{iq:rs:portfolio-construction-i:2}
What is a shadow price in portfolio construction, and how would you use it?
\end{interviewq}

\begin{interviewq}[$\star\star$ \iqroles{researcher, risk}]\label{iq:rs:portfolio-construction-i:3}
Your book must be dollar, beta and sector neutral. What does that cost, and when is it free?
\end{interviewq}

\begin{interviewq}[$\star\star$ \iqroles{trader}]\label{iq:rs:portfolio-construction-i:4}
How do you set a liquidity limit, and what happens to the book as capital grows?
\end{interviewq}

\begin{interviewq}[$\star\star$ \iqroles{researcher}]\label{iq:rs:portfolio-construction-i:5}
What is the \omterm{def:rs:from-signal-to-forecast:breadth}{transfer coefficient}, and why can a lower one be better?
\end{interviewq}

\begin{interviewq}[$\star\star\star$ \iqroles{researcher, developer}]\label{iq:rs:portfolio-construction-i:6}
Design the optimiser for a book of three thousand names rebalanced daily: formulation, solver, and what you precompute.
\end{interviewq}
